\title{Επισκόπηση Λειτουργίας Συστήματος Μεταφοράς Ηλεκτρικής Ενέργειας}

\section{Θεωρία}
\label{kef-01-theoria}

Το κεφάλαιο αυτό πραγματεύεται ένα ερώτημα που φαίνεται απλό χωρίς να είναι: με ποιον
τρόπο διατηρείται σε ασφαλή λειτουργία ένα σύστημα ηλεκτρικής ενέργειας. Η απάντηση
συνίσταται σε έναν κλειστό βρόχο τεσσάρων σταδίων — μέτρηση του συστήματος, μετατροπή των
μετρήσεων σε αξιόπιστη εικόνα, αξιολόγηση της ασφάλειας της εικόνας αυτής, και δράση. Οι
ενότητες που ακολουθούν εξετάζουν διαδοχικά τα στάδια αυτά. Το ζητούμενο δεν είναι η
απομνημόνευση των ακρωνυμίων, αλλά η κατανόηση του λόγου ύπαρξης κάθε σταδίου και των
συνεπειών της απουσίας του.

\subsection{Ο ρόλος του Συστήματος Μεταφοράς}

Το Σύστημα Μεταφοράς Ηλεκτρικής Ενέργειας (ΣΜΗΕ) είναι το δίκτυο υπερυψηλής και υψηλής
τάσης που διασυνδέει τους σταθμούς παραγωγής με τα κέντρα κατανάλωσης και με τα γειτονικά
συστήματα. Στην Ελλάδα λειτουργεί στα επίπεδα των 400 kV, 150 kV και 66 kV, ενώ τη
διαχείρισή του ασκεί ο Ανεξάρτητος Διαχειριστής Μεταφοράς Ηλεκτρικής Ενέργειας (ΑΔΜΗΕ), ο
οποίος αποτελεί μέλος του ENTSO-E και τμήμα της σύγχρονης περιοχής της Ηπειρωτικής Ευρώπης
(Continental Europe).

Ο λόγος για τον οποίο η μεταφορά γίνεται σε τόσο υψηλή τάση προκύπτει άμεσα από τις
θεμελιώδεις σχέσεις. Για δεδομένη μεταφερόμενη ισχύ $P = \sqrt{3}\,V I \cos\varphi$, το
ρεύμα είναι αντιστρόφως ανάλογο της τάσης, ενώ οι απώλειες στους αγωγούς είναι ανάλογες
του $I^2$. Διπλασιασμός της τάσης υποτετραπλασιάζει επομένως τις απώλειες μεταφοράς. Σε
αυτό οφείλεται και η ιεραρχική οργάνωση του συστήματος: υπερυψηλή τάση για τις μεγάλες
αποστάσεις και διαδοχικοί υποβιβασμοί προς την πλευρά του καταναλωτή.

\subsubsection*{Οι τρεις ταυτόχρονες απαιτήσεις}

Ο Διαχειριστής Συστήματος Μεταφοράς (Transmission System Operator, TSO) καλείται να
ικανοποιεί ταυτόχρονα τρεις απαιτήσεις. Η μεταξύ τους ένταση είναι εκείνη που προσδίδει
στο πρόβλημα τη δυσκολία του:

\begin{itemize}
\item \textbf{Ισοζύγιο ισχύος.} Η ηλεκτρική ενέργεια δεν αποθηκεύεται σε σημαντικές
      ποσότητες εντός του συστήματος. Παραγωγή και ζήτηση οφείλουν να εξισώνονται
      \emph{συνεχώς}. Σε αντίθετη περίπτωση, η διαφορά καλύπτεται προσωρινά από την
      κινητική ενέργεια των στρεφόμενων μαζών, με τίμημα τη μεταβολή της συχνότητας: ένα
      σύστημα που καταναλώνει περισσότερο από όσο παράγει επιβραδύνεται κυριολεκτικά.
\item \textbf{Ασφάλεια λειτουργίας.} Το σύστημα οφείλει να παραμένει εντός των ορίων
      ασφαλούς λειτουργίας — θερμικών ορίων στοιχείων, ορίων τάσεων, ορίων ευστάθειας —
      όχι μόνο στην τρέχουσα κατάσταση, αλλά και μετά από κάθε πιθανή μεμονωμένη διαταραχή.
      Η διάκριση αυτή μεταξύ της τρέχουσας και της μετά τη διαταραχή κατάστασης είναι
      θεμελιώδης και οδηγεί στο κριτήριο N-1, το οποίο ορίζεται αμέσως παρακάτω.
\item \textbf{Ποιότητα.} Η συχνότητα και οι τάσεις οφείλουν να διατηρούνται εντός
      προδιαγεγραμμένων περιθωρίων, τα οποία για την Ηπειρωτική Ευρώπη ορίζονται
      ρυθμιστικά και όχι κατά την κρίση του Διαχειριστή.
\end{itemize}

Η δυσκολία έγκειται στο ότι οι απαιτήσεις αυτές εκδηλώνονται σε \emph{διαφορετικές χρονικές
κλίμακες}, οι οποίες εκτείνονται σε πάνω από δώδεκα τάξεις μεγέθους. Ο
Πίνακας~\ref{tab:chronikes-klimakes} συνοψίζει την εικόνα και λειτουργεί ως χάρτης για το
υπόλοιπο κεφάλαιο.

\begin{table}[!htbp]
\centering
\begin{tabular}{lll}
\hline
Κλίμακα & Φαινόμενο ή ενέργεια & Μηχανισμός \\
\hline
μs -- ms   & κύματα, σφάλματα           & προστασία, διακόπτες \\
ms -- s    & ηλεκτρομηχανικές ταλαντώσεις & αδράνεια, PSS, FACTS \\
s -- 30 s  & πτώση συχνότητας            & εφεδρεία διατήρησης συχνότητας (FCR) \\
30 s -- 15 min & αποκατάσταση συχνότητας & εφεδρεία αποκατάστασης (aFRR, mFRR) \\
min -- h   & αναπρογραμματισμός          & αγορά εξισορρόπησης, ενδοημερήσια \\
h -- ημέρα & κατανομή φορτίου            & αγορά επόμενης ημέρας \\
εβδομάδες -- έτη & επάρκεια, συντηρήσεις & προγραμματισμός, επενδύσεις \\
\hline
\end{tabular}
\caption{Χρονικές κλίμακες λειτουργίας του ΣΜΗΕ. Σε κάθε γραμμή αντιστοιχεί διαφορετικός
μηχανισμός ελέγχου· το παρόν κεφάλαιο εστιάζει από τα δευτερόλεπτα και άνω.}
\label{tab:chronikes-klimakes}
\end{table}

\subsection{Καταστάσεις λειτουργίας και κριτήριο N-1}

Ένα σύστημα δεν χαρακτηρίζεται απλώς ως ασφαλές ή μη ασφαλές. Η καθιερωμένη σήμερα εικόνα
οφείλεται στον Dy Liacco, ο οποίος στα τέλη της δεκαετίας του 1960 διατύπωσε την αρχή ότι
ο έλεγχος οφείλει να \emph{προσαρμόζεται} στην εκάστοτε κατάσταση ασφαλείας. Το ευρωπαϊκό
κανονιστικό πλαίσιο (System Operation Guideline, Κανονισμός ΕΕ 2017/1485, Άρθρο 18)
κωδικοποιεί την αρχή αυτή σε πέντε καταστάσεις:

\begin{enumerate}
\item \textbf{Κανονική} (normal). Όλα τα μεγέθη βρίσκονται εντός ορίων \emph{και} το
      κριτήριο N-1 ικανοποιείται. Ο έλεγχος είναι \emph{προληπτικός}: ο Διαχειριστής
      προετοιμάζεται για διαταραχές που δεν έχουν εκδηλωθεί.
\item \textbf{Επιφυλακής} (alert). Τα μεγέθη παραμένουν εντός ορίων, έχει όμως εξαντληθεί
      το περιθώριο εφεδρείας ή παραβιάζεται το κριτήριο N-1. Το σύστημα λειτουργεί, αλλά
      μια επόμενη διαταραχή θα το οδηγήσει εκτός ορίων.
\item \textbf{Έκτακτης ανάγκης} (emergency). Παραβιάζεται τουλάχιστον ένα όριο ασφαλούς
      λειτουργίας. Ο έλεγχος καθίσταται \emph{διορθωτικός} και επείγων.
\item \textbf{Συσκότισης} (blackout). Εκτεταμένη απώλεια φορτίου ή κατάρρευση τμήματος του
      συστήματος.
\item \textbf{Αποκατάστασης} (restoration). Ενεργοποιείται το σχέδιο επανατροφοδότησης.
\end{enumerate}

Η μετάβαση από την κανονική στην κατάσταση επιφυλακής \emph{δεν} συνοδεύεται από ορατή
μεταβολή των μετρούμενων μεγεθών, καθώς κανένα όριο δεν υπερβαίνεται. Πρόκειται για
μεταβολή που υφίσταται μόνο στους υπολογισμούς του Διαχειριστή. Ακριβώς γι' αυτό απαιτείται
αξιόπιστη εικόνα του συστήματος: χωρίς αυτήν, η διαφορά μεταξύ ασφαλούς και επικίνδυνης
λειτουργίας παραμένει αόρατη.

\subsubsection*{Το κριτήριο N-1}

Το \textbf{κριτήριο N-1} αποτελεί τον πρακτικό ορισμό της ασφάλειας: από τα $N$ στοιχεία
του συστήματος, η απώλεια οποιουδήποτε \emph{ενός} — γραμμής, μετασχηματιστή, μονάδας
παραγωγής ή ζυγού — πρέπει να μπορεί να απορροφηθεί χωρίς παραβίαση ορίων και χωρίς
αλυσιδωτές αποσυνδέσεις.

Ο έλεγχός του πραγματοποιείται με \emph{ανάλυση ενδεχομένων} (contingency analysis): για
κάθε πιθανό ενδεχόμενο εκτελείται ροή ισχύος με το αντίστοιχο στοιχείο εκτός λειτουργίας
και ελέγχεται τυχόν υπέρβαση ορίων. Σε σύστημα με μερικές εκατοντάδες κλάδους, αυτό
συνεπάγεται εκατοντάδες ροές ισχύος που πρέπει να ολοκληρωθούν εντός ολίγων λεπτών.

Δύο παρατηρήσεις αξίζει να συγκρατηθούν:

\begin{itemize}
\item Το N-1 αποτελεί \emph{ντετερμινιστικό} κριτήριο, καθώς δεν σταθμίζει την πιθανότητα
      κάθε ενδεχομένου. Ένα σπάνιο και ένα συχνό σφάλμα αντιμετωπίζονται ισοδύναμα. Είναι
      απλό και ελέγξιμο, αλλά συντηρητικό σε ορισμένες περιπτώσεις και ανεπαρκές σε άλλες —
      εξ ου και το ενδιαφέρον για πιθανοτικά κριτήρια ασφαλείας.
\item Η ανάλυση ενδεχομένων απαιτεί ως \emph{είσοδο} την τρέχουσα κατάσταση του συστήματος.
      Αν η είσοδος αυτή είναι εσφαλμένη, εσφαλμένα είναι και όλα τα συμπεράσματα περί
      ασφάλειας. Το γεγονός αυτό καθιστά την εκτίμηση κατάστασης, η οποία εξετάζεται
      παρακάτω, όχι απλώς χρήσιμη αλλά προαπαιτούμενη.
\end{itemize}

\subsection{Εποπτεία και έλεγχος: SCADA και EMS}

\subsubsection*{Η αλυσίδα μέτρησης}

Η βασική υποδομή τηλεμέτρησης και τηλεχειρισμού του ΣΜΗΕ είναι το σύστημα \textbf{SCADA}
(Supervisory Control and Data Acquisition). Η αλυσίδα από το φυσικό μέγεθος έως την οθόνη
του χειριστή περιλαμβάνει τέσσερις κρίκους:

\begin{enumerate}
\item \textbf{Μετασχηματιστές οργάνων.} Οι μετασχηματιστές τάσης (VT) και έντασης (CT)
      υποβιβάζουν τα μεγέθη του δικτύου σε επίπεδα διαχειρίσιμα από τα ηλεκτρονικά
      κυκλώματα. Εισάγουν το πρώτο σφάλμα, τυπικά κλάσης 0,2 έως 1\%.
\item \textbf{Μεταλλάκτες και ευφυείς ηλεκτρονικές συσκευές (IED).} Υπολογίζουν τα παράγωγα
      μεγέθη — ενεργό και άεργο ισχύ, ενεργό τιμή τάσης — από τα στιγμιαία δείγματα.
\item \textbf{Απομακρυσμένες τερματικές μονάδες (RTU).} Συγκεντρώνουν τα δεδομένα του
      υποσταθμού και τα αποστέλλουν στο Κέντρο Ελέγχου Ενέργειας.
\item \textbf{Τηλεπικοινωνιακό δίκτυο.} Μεταφέρει τα δεδομένα με τυποποιημένα πρωτόκολλα,
      με πλέον διαδεδομένα τα IEC 60870-5-101/104 και το DNP3, ενώ εντός του υποσταθμού
      κυριαρχεί πλέον το IEC 61850.
\end{enumerate}

Συλλέγονται δύο κατηγορίες δεδομένων:

\begin{itemize}
\item \textbf{Αναλογικές μετρήσεις}: ενεργός και άεργος ροή ισχύος στα άκρα των γραμμών και
      των μετασχηματιστών, εγχεόμενη ισχύς σε ζυγούς, μέτρα τάσεων, ρεύματα κλάδων.
\item \textbf{Ψηφιακές ενδείξεις}: θέσεις διακοπτών και αποζευκτών, θέσεις λήψεων
      μετασχηματιστών, σήματα συναγερμών και λειτουργίας προστασιών. Οι ενδείξεις αυτές
      καθορίζουν την \emph{τοπολογία}, δηλαδή ποια στοιχεία βρίσκονται εν λειτουργία και
      με ποιον τρόπο συνδέονται μεταξύ τους.
\end{itemize}

Ο ρυθμός σάρωσης συμβατικού συστήματος SCADA είναι τυπικά της τάξης ολίγων δευτερολέπτων,
συνήθως 2 έως 10 s.

\subsubsection*{Δύο καθοριστικοί περιορισμοί}

Δύο χαρακτηριστικά του SCADA είναι κρίσιμα για όσα ακολουθούν και αξίζει να συγκρατηθούν:

\begin{itemize}
\item Οι μετρήσεις \textbf{δεν είναι χρονικά συγχρονισμένες} μεταξύ τους. Συλλέγονται σε
      διαφορετικές χρονικές στιγμές εντός του κύκλου σάρωσης, οπότε το σύνολό τους αποτελεί
      \emph{οιονεί στατική} απεικόνιση και όχι στιγμιότυπο. Ένα σύνολο μετρήσεων που έχει
      κατανεμηθεί σε πέντε δευτερόλεπτα δεν αντιστοιχεί σε καμία πραγματική στιγμή του
      συστήματος. Το γεγονός αυτό είναι αμελητέο όσο το σύστημα μεταβάλλεται αργά, αλλά
      καθίσταται σοβαρό κατά τη διάρκεια διαταραχής.
\item Οι μετρήσεις είναι \textbf{θορυβώδεις και ενίοτε εσφαλμένες}, ενώ ορισμένα μεγέθη δεν
      μετρώνται καθόλου. Χαρακτηριστικά, οι \emph{γωνίες} των τάσεων — οι οποίες αποτελούν
      το ήμισυ της κατάστασης του συστήματος — απουσιάζουν εξ ολοκλήρου από το συμβατικό
      SCADA.
\end{itemize}

Ένας απορυθμισμένος μετασχηματιστής έντασης, μια RTU με εσφαλμένη πολικότητα, ένας
διακόπτης του οποίου η ένδειξη θέσης έχει παγώσει λόγω βλάβης: όλες αυτές οι περιπτώσεις
παράγουν δεδομένα που εμφανίζονται απολύτως φυσιολογικά. Το σύστημα εποπτείας δεν διαθέτει
τρόπο να τα διακρίνει \emph{ανά μέτρηση}. Η διάκριση είναι εφικτή μόνο συλλογικά, με
αξιοποίηση του πλεονασμού, όπως αναλύεται στη συνέχεια.

\subsubsection*{Το EMS και η σειρά των λειτουργιών}

Το \textbf{EMS} (Energy Management System) είναι το λογισμικό που μετατρέπει τα ακατέργαστα
δεδομένα σε αποφάσεις. Η σειρά εκτέλεσης των λειτουργιών του δεν είναι αυθαίρετη, καθώς
κάθε βήμα τροφοδοτεί το επόμενο:

\begin{enumerate}
\item \textbf{Επεξεργαστής τοπολογίας} (topology processor). Συνθέτει, από τις θέσεις των
      διακοπτών, το ηλεκτρικό μοντέλο ζυγών-κλάδων (bus-branch model). Προσδιορίζει τον
      τρόπο με τον οποίο είναι συνδεδεμένο το δίκτυο τη δεδομένη στιγμή.
\item \textbf{Εκτίμηση κατάστασης}. Προσδιορίζει την ηλεκτρική κατάσταση του δικτύου αυτού.
\item \textbf{Ανάλυση ενδεχομένων}. Αποφαίνεται για την ασφάλεια της κατάστασης αυτής.
\item \textbf{Βέλτιστη ροή ισχύος} (Optimal Power Flow). Προσδιορίζει, εφόσον απαιτείται,
      τις αναγκαίες διορθωτικές ενέργειες και το κόστος τους.
\item \textbf{Αυτόματος έλεγχος παραγωγής} (Automatic Generation Control). Υλοποιεί συνεχώς
      τη ρύθμιση συχνότητας και τις διασυνοριακές ανταλλαγές.
\end{enumerate}

Σφάλμα στο πρώτο ή στο δεύτερο βήμα διαδίδεται σιωπηρά σε όλα τα επόμενα.

Το Σχήμα~\ref{fig:scada-ems} συνοψίζει την πλήρη διαδρομή. Αξίζει να παρατηρηθεί ότι
πρόκειται για \emph{κλειστό} βρόχο: τα δεδομένα ανεβαίνουν από το δίκτυο προς το κέντρο
ελέγχου και οι εντολές επιστρέφουν σε αυτό.

\begin{figure}[!htbp]
\centering
\includegraphics[width=\linewidth]{figures/scada-ems.svg}
\caption{Ο βρόχος εποπτείας και ελέγχου. Η επάνω σειρά είναι η αλυσίδα μέτρησης, από το
φυσικό δίκτυο έως το SCADA· η κάτω σειρά είναι η ακολουθία λειτουργιών του EMS, η οποία
καταλήγει σε εντολές προς το δίκτυο.}
\label{fig:scada-ems}
\end{figure}

\subsection{Συγχρονισμένες μετρήσεις φασιθετών και συστήματα WAMS}

\subsubsection*{Η σημασία της γωνίας}

Υπενθυμίζεται η σχέση μεταφοράς ενεργού ισχύος μεταξύ δύο ζυγών που συνδέονται μέσω
επαγωγικής γραμμής αντίδρασης $X$:

\begin{equation}
P_{12} = \frac{V_1 V_2}{X}\,\sin(\delta_1 - \delta_2)
\label{eq:powerflow}
\end{equation}

Η ροή ισχύος καθορίζεται από τη \emph{διαφορά γωνιών}. Η γωνία δεν αποτελεί επομένως
βοηθητικό μαθηματικό μέγεθος, αλλά την κινητήρια αιτία της ροής ισχύος και τον πλέον άμεσο
δείκτη καταπόνησης του δικτύου. Η αδυναμία μέτρησής της από το συμβατικό SCADA συνιστά
σοβαρή έλλειψη.

\subsubsection*{Η αρχή λειτουργίας}

Οι \textbf{μονάδες μέτρησης φασιθετών} (Phasor Measurement Units, PMU) αίρουν τον
περιορισμό αυτόν. Μια PMU υπολογίζει τον φασιθέτη — μέτρο \emph{και} γωνία — της τάσης και
του ρεύματος, και τον συνοδεύει με χρονοσήμανση από κοινή χρονική αναφορά UTC, τυπικά μέσω
GPS.

Το ουσιώδες σημείο έγκειται στο εξής: η γωνία ενός φασιθέτη έχει νόημα μόνο ως προς κάποια
αναφορά. Επειδή όλες οι PMU ανάγονται στην ίδια χρονική βάση, οι γωνίες που μετρούν είναι
άμεσα συγκρίσιμες σε ολόκληρο το σύστημα, ακόμη και σε αποστάσεις εκατοντάδων χιλιομέτρων —
εξ ου και ο όρος \emph{συγχρονοφασιθέτες} (synchrophasors). Για να είναι έγκυρη η σύγκριση
σε σύστημα 50 Hz απαιτείται ακρίβεια χρονισμού καλύτερη του μικροδευτερολέπτου: σφάλμα 1 μs
αντιστοιχεί σε σφάλμα γωνίας $360^\circ \times 50 \times 10^{-6} = 0{,}018^\circ$.

\begin{figure}[!htbp]
\centering
\includegraphics[width=\linewidth]{figures/pmu-wams.svg}
\caption{Αριστερά, το μετρούμενο μέγεθος: ο φασιθέτης τάσης με μέτρο και γωνία ως προς
κοινή χρονική αναφορά. Δεξιά, η ιεραρχία συλλογής: οι PMU συγχρονίζονται μέσω GPS, τα
πλαίσιά τους ευθυγραμμίζονται χρονικά στον PDC και τροφοδοτούν το WAMS.}
\label{fig:pmu-wams}
\end{figure}

\subsubsection*{Προδιαγραφές}

Τα βασικά τεχνικά χαρακτηριστικά καθορίζονται από το πρότυπο IEC/IEEE 60255-118-1:2018, το
οποίο διαδέχθηκε τη σειρά IEEE C37.118:

\begin{itemize}
\item \textbf{Ρυθμός αναφοράς}: τυπικά 10, 25 ή 50 πλαίσια ανά δευτερόλεπτο για σύστημα
      50 Hz, με δυνατότητα υψηλότερων ρυθμών.
\item \textbf{Ακρίβεια}: το συνολικό σφάλμα φασιθέτη (Total Vector Error, TVE) δεν
      επιτρέπεται να υπερβαίνει το 1\% στη μόνιμη κατάσταση. Ορίζεται ως το μέτρο της
      διαφοράς του μετρούμενου από τον πραγματικό φασιθέτη, κανονικοποιημένο ως προς τον
      πραγματικό, και συνεπώς ενοποιεί σε ένα μέγεθος τα σφάλματα μέτρου και γωνίας.
      Ορίζονται επιπλέον το σφάλμα συχνότητας (FE) και το σφάλμα ρυθμού μεταβολής
      συχνότητας (RFE).
\item \textbf{Κλάσεις}: η κλάση P (protection) βελτιστοποιείται για ταχεία απόκριση, με
      μικρή καθυστέρηση φίλτρου, ενώ η κλάση M (measurement) για ακρίβεια υπό την παρουσία
      αρμονικών και παρεμβολών. Η επιλογή συνιστά συμβιβασμό, καθώς δεν υφίσταται φίλτρο
      ταυτοχρόνως ταχύ και ισχυρά απορριπτικό.
\end{itemize}

\begin{table}[!htbp]
\centering
\begin{tabular}{lll}
\hline
 & SCADA & PMU / WAMS \\
\hline
Ρυθμός           & 2--10 s ανά σάρωση & 10--50 πλαίσια ανά s \\
Συγχρονισμός     & όχι                & ναι (UTC μέσω GPS) \\
Γωνία τάσης      & δεν μετράται       & μετράται άμεσα \\
Φύση μετρήσεων   & μη γραμμικές ως προς την κατάσταση & γραμμικές \\
Κάλυψη           & εκτεταμένη         & επιλεκτική, λόγω κόστους \\
Τυπική χρήση     & εκτίμηση κατάστασης, τηλεχειρισμός & δυναμικά φαινόμενα \\
\hline
\end{tabular}
\caption{Σύγκριση συμβατικής και συγχρονισμένης μέτρησης. Οι δύο τεχνολογίες είναι
συμπληρωματικές και όχι ανταγωνιστικές.}
\label{tab:scada-pmu}
\end{table}

\subsubsection*{Από τη μέτρηση στην εφαρμογή}

Οι μετρήσεις συγκεντρώνονται ιεραρχικά σε συμπυκνωτές δεδομένων φασιθετών (Phasor Data
Concentrators, PDC), οι οποίοι ευθυγραμμίζουν χρονικά τα πλαίσια από διαφορετικές PMU και
συνθέτουν ένα \textbf{σύστημα εποπτείας ευρείας περιοχής} (Wide Area Monitoring System,
WAMS). Οι χαρακτηριστικές εφαρμογές είναι:

\begin{itemize}
\item η παρακολούθηση ηλεκτρομηχανικών ταλαντώσεων μεταξύ περιοχών, με συχνότητες τυπικά
      0,1 έως 1 Hz, οι οποίες παραμένουν πρακτικά αόρατες σε ρυθμό σάρωσης δευτερολέπτων·
\item η μέτρηση του ρυθμού μεταβολής της συχνότητας και η εκτίμηση της διαθέσιμης
      αδράνειας·
\item η ανίχνευση νησιδοποίησης και η επιτήρηση της γωνιακής απόκλισης μεταξύ περιοχών·
\item η επικύρωση μοντέλων μετά από διαταραχές, μέσω σύγκρισης της καταγεγραμμένης με την
      προσομοιωμένη απόκριση·
\item η ενίσχυση — ή, σε συνθήκες πλήρους κάλυψης, η αντικατάσταση — της κλασικής εκτίμησης
      κατάστασης.
\end{itemize}

\subsection{Ευέλικτα συστήματα μεταφοράς εναλλασσομένου ρεύματος (FACTS)}

Ενώ τα συστήματα SCADA και PMU \emph{παρατηρούν} το σύστημα, οι διατάξεις \textbf{FACTS}
(Flexible AC Transmission Systems) το \emph{ελέγχουν} ενεργά.

Η \eqref{eq:powerflow} προσδιορίζει με ακρίβεια τα ελεγχόμενα μεγέθη. Η ροή ισχύος
εξαρτάται από τρία και μόνο μεγέθη: τις τάσεις $V_1, V_2$, την αντίδραση $X$ και τη διαφορά
γωνιών $\delta_1 - \delta_2$. Σε κλασικό δίκτυο εναλλασσομένου ρεύματος και τα τρία είναι
ουσιαστικά δεδομένα, οπότε η ισχύς ρέει σύμφωνα με τους νόμους της φυσικής και όχι σύμφωνα
με τις επιδιώξεις του Διαχειριστή. Οι διατάξεις FACTS είναι μετατροπείς ηλεκτρονικών ισχύος
που καθιστούν τα μεγέθη αυτά ελεγχόμενα, ταχέως και συνεχώς. Κατατάσσονται ανάλογα με το
μέγεθος που επηρεάζουν, όπως συνοψίζει το Σχήμα~\ref{fig:facts}:

\begin{itemize}
\item \textbf{Παράλληλης σύνδεσης} (έλεγχος του $V$): ο στατός αντισταθμιστής αέργου ισχύος
      (Static Var Compensator, SVC) και ο STATCOM. Ελέγχουν την τάση του ζυγού μέσω έγχυσης
      ή απορρόφησης αέργου ισχύος.
\item \textbf{Σειριακής σύνδεσης} (έλεγχος του $X$): ο ελεγχόμενος με θυρίστορ πυκνωτής
      σειράς (Thyristor Controlled Series Capacitor, TCSC) και ο SSSC. Μεταβάλλοντας τη
      φαινόμενη αντίδραση της γραμμής, ανακατευθύνουν τη ροή ισχύος και αυξάνουν τη
      μεταφορική ικανότητα.
\item \textbf{Συνδυασμένης λειτουργίας} (έλεγχος και της γωνίας): ο ενοποιημένος ελεγκτής
      ροής ισχύος (Unified Power Flow Controller, UPFC) ελέγχει ταυτόχρονα και ανεξάρτητα
      τάση, ενεργό και άεργο ροή. Αποτελεί την πληρέστερη και ταυτόχρονα την ακριβότερη
      λύση.
\end{itemize}

\begin{figure}[!htbp]
\centering
\includegraphics[width=0.88\linewidth]{figures/facts.svg}
\caption{Οι τρεις κατηγορίες διατάξεων FACTS σε γραμμή δύο ζυγών. Κάθε κατηγορία επεμβαίνει
σε διαφορετικό μέγεθος της σχέσης μεταφοράς ισχύος: η παράλληλη στην τάση, η σειριακή στην
αντίδραση, η συνδυασμένη και στη γωνία.}
\label{fig:facts}
\end{figure}

Σε συνδυασμό με τους συνδέσμους συνεχούς ρεύματος υψηλής τάσης (HVDC), όπου η ροή ισχύος
ορίζεται απευθείας από τον έλεγχο του μετατροπέα, οι διατάξεις FACTS αποτελούν τα βασικά
εργαλεία του Διαχειριστή για τη ρύθμιση τάσης και την ανακούφιση συμφορήσεων στο επίπεδο
μεταφοράς. Τα αντίστοιχα εργαλεία στο επίπεδο διανομής εξετάζονται στα Κεφάλαια 5 και 6.

\subsection{Εκτίμηση κατάστασης}

\subsubsection*{Διατύπωση του προβλήματος}

Όπως αναλύθηκε παραπάνω, οι μετρήσεις του SCADA είναι πλεονάζουσες, θορυβώδεις, ασύγχρονες
και ελλιπείς, ενώ παράλληλα η ανάλυση ασφάλειας απαιτεί \emph{ενιαία και συνεπή} εικόνα του
δικτύου. Η \textbf{εκτίμηση κατάστασης} (state estimation) γεφυρώνει το χάσμα αυτό.
Εισήχθη στα συστήματα ηλεκτρικής ενέργειας από τους Schweppe και Wildes το 1970, σε σειρά
τριών εργασιών οι οποίες παραμένουν έως σήμερα σημείο αναφοράς, και αποτελεί έκτοτε τη
θεμελιώδη λειτουργία κάθε EMS.

Μια πρώτη προσέγγιση θα ήταν η απλή επίλυση μιας ροής ισχύος. Η προσέγγιση αυτή δεν
επαρκεί, για δύο λόγους. Πρώτον, η ροή ισχύος απαιτεί \emph{ακριβώς} τα κατάλληλα δεδομένα
εισόδου, ενώ τα διαθέσιμα δεδομένα είναι πολύ περισσότερα από τα αναγκαία και όλα ελαφρώς
εσφαλμένα. Δεύτερον, η ροή ισχύος δεν διαθέτει μηχανισμό ανίχνευσης εσφαλμένων δεδομένων
και θα συγκλίνει απρόσκοπτα σε εσφαλμένη λύση. Απαιτείται επομένως μέθοδος που
\emph{αξιοποιεί} την πλεονάζουσα πληροφορία αντί να την αγνοεί.

Το Σχήμα~\ref{fig:se-flow} δίνει τη συνολική ροή του αλγορίθμου. Οι επιμέρους βαθμίδες
αναλύονται στις υποενότητες που ακολουθούν· χρήσιμο είναι να σημειωθεί από τώρα ότι
υπάρχουν \emph{δύο} εμφωλευμένοι βρόχοι — ο εσωτερικός για τη σύγκλιση της αριθμητικής
επίλυσης και ο εξωτερικός για την απόρριψη εσφαλμένων μετρήσεων.

\begin{figure}[!htbp]
\centering
\includegraphics[width=0.8\linewidth]{figures/se-flowchart.svg}
\caption{Ροή της εκτίμησης κατάστασης σε ένα EMS. Ο εσωτερικός βρόχος (αριστερά) επαναλαμβάνει
τη μέθοδο Gauss--Newton έως τη σύγκλιση· ο εξωτερικός (δεξιά) απορρίπτει τη μέτρηση με το
μέγιστο κανονικοποιημένο υπόλοιπο και επανεκτιμά.}
\label{fig:se-flow}
\end{figure}

\subsubsection*{Το μοντέλο μέτρησης}

Ως \emph{κατάσταση} ορίζεται το διάνυσμα
$\mathbf{x} = [\boldsymbol{\theta}^{T}, \mathbf{V}^{T}]^{T}$ των γωνιών και των μέτρων των
τάσεων όλων των ζυγών. Για δίκτυο $N$ ζυγών, η γωνία ενός ζυγού ορίζεται αυθαίρετα ως
αναφορά, συνήθως μηδενική, οπότε το πλήθος των αγνώστων ανέρχεται σε

\begin{equation}
n = 2N - 1
\label{eq:nstates}
\end{equation}

Η επιλογή αυτή είναι θεμελιώδης: εφόσον είναι γνωστά τα $\mathbf{V}$ και
$\boldsymbol{\theta}$, \emph{κάθε} άλλο μέγεθος του δικτύου — ροές, εγχύσεις, απώλειες —
υπολογίζεται από τις εξισώσεις δικτύου. Η κατάσταση αποτελεί το ελάχιστο σύνολο μεγεθών που
περιγράφει πλήρως το σύστημα.

Το μοντέλο μέτρησης δίνεται από τη σχέση

\begin{equation}
\mathbf{z} = \mathbf{h}(\mathbf{x}) + \mathbf{e}
\label{eq:measmodel}
\end{equation}

όπου $\mathbf{z} \in \mathbb{R}^{m}$ το διάνυσμα των μετρήσεων, $\mathbf{h}(\cdot)$ οι μη
γραμμικές συναρτήσεις που συνδέουν την κατάσταση με τα μετρούμενα μεγέθη, και $\mathbf{e}$
ο θόρυβος μέτρησης. Ο τελευταίος θεωρείται μηδενικής μέσης τιμής, με διαγώνιο πίνακα
συνδιακύμανσης $\mathbf{R} = \mathrm{diag}(\sigma_1^2, \ldots, \sigma_m^2)$, δηλαδή τα
σφάλματα των επιμέρους μετρητών θεωρούνται ασυσχέτιστα.

Καθοριστικό μέγεθος αποτελεί ο \textbf{λόγος πλεονασμού}

\begin{equation}
\eta = \frac{m}{n}
\label{eq:redundancy}
\end{equation}

ο οποίος σε πραγματικά συστήματα κυμαίνεται τυπικά μεταξύ 1,7 και 3. Ενδεικτικά, σε δίκτυο
$N = 100$ ζυγών οι άγνωστοι ανέρχονται σε $n = 199$ και, για $\eta = 2$, οι μετρήσεις σε
$m = 400$ περίπου. Ο πλεονασμός είναι εκείνος που επιτρέπει τόσο τη μείωση του θορύβου όσο
και, το κυριότερο, την ανίχνευση εσφαλμένων μετρήσεων. Ελλείψει αυτού, το πρόβλημα
εκφυλίζεται σε ροή ισχύος, με όλα τα μειονεκτήματα που προαναφέρθηκαν.

\subsubsection*{Σταθμισμένα ελάχιστα τετράγωνα}

Επειδή οι μετρήσεις είναι περισσότερες από τους αγνώστους, δεν υφίσταται $\mathbf{x}$ που
να τις ικανοποιεί όλες. Αναζητείται επομένως εκείνο που τις ικανοποιεί κατά τον βέλτιστο
δυνατό τρόπο. Η καθιερωμένη διατύπωση είναι εκείνη των \textbf{σταθμισμένων ελαχίστων
τετραγώνων} (Weighted Least Squares, WLS):

\begin{equation}
J(\mathbf{x}) = \sum_{i=1}^{m} \frac{\left(z_i - h_i(\mathbf{x})\right)^2}{\sigma_i^2}
             = \left[\mathbf{z}-\mathbf{h}(\mathbf{x})\right]^{T} \mathbf{R}^{-1}
               \left[\mathbf{z}-\mathbf{h}(\mathbf{x})\right]
\label{eq:wls}
\end{equation}

Η στάθμιση με $1/\sigma_i^2$ έχει άμεση φυσική σημασία: μια μέτρηση με μικρή τυπική
απόκλιση, δηλαδή ακριβής, βαρύνει περισσότερο, ενώ μια θορυβώδης βαρύνει λιγότερο.
Ισοδύναμα, κάθε όρος του αθροίσματος αποτελεί το τετράγωνο του σφάλματος \emph{εκφρασμένο
σε μονάδες της οικείας τυπικής απόκλισης}, ώστε να καθίστανται συγκρίσιμα μεταξύ τους
μεγέθη διαφορετικής φύσης και κλίμακας. Υπό την υπόθεση κανονικά κατανεμημένων σφαλμάτων,
η ελαχιστοποίηση της \eqref{eq:wls} ταυτίζεται με την εκτίμηση μέγιστης πιθανοφάνειας.

Επειδή η $\mathbf{h}$ είναι μη γραμμική, η λύση προκύπτει επαναληπτικά με τη μέθοδο
Gauss--Newton, μέσω διαδοχικών βημάτων γραμμικοποίησης γύρω από την τρέχουσα εκτίμηση,
επίλυσης του γραμμικού προβλήματος και ενημέρωσης. Με
$\mathbf{H}(\mathbf{x}) = \partial \mathbf{h}/\partial \mathbf{x}$ τον ιακωβιανό πίνακα, σε
κάθε επανάληψη $k$ επιλύεται το σύστημα των \emph{κανονικών εξισώσεων}

\begin{align}
\mathbf{G}(\mathbf{x}^{k}) \, \Delta \mathbf{x}^{k}
  &= \mathbf{H}^{T}(\mathbf{x}^{k}) \, \mathbf{R}^{-1}
     \left[\mathbf{z} - \mathbf{h}(\mathbf{x}^{k})\right] \\
\mathbf{G}(\mathbf{x}^{k})
  &= \mathbf{H}^{T}(\mathbf{x}^{k}) \, \mathbf{R}^{-1} \mathbf{H}(\mathbf{x}^{k})
\label{eq:gain}
\end{align}

και ενημερώνεται η εκτίμηση ως $\mathbf{x}^{k+1} = \mathbf{x}^{k} + \Delta \mathbf{x}^{k}$,
έως ότου το $\max|\Delta \mathbf{x}^{k}|$ κατέλθει κάτω από προκαθορισμένο κατώφλι
σύγκλισης. Η σύγκλιση είναι τυπικά ταχεία, της τάξης των τριών έως πέντε επαναλήψεων,
επειδή ως αρχική εκτίμηση χρησιμοποιείται συνήθως η λύση της προηγούμενης εκτέλεσης, ενώ το
σύστημα δεν έχει μεταβληθεί σημαντικά στο μεσοδιάστημα.

Ο πίνακας $\mathbf{G}$ ονομάζεται \emph{πίνακας κέρδους} (gain matrix). Είναι συμμετρικός,
θετικά ημιορισμένος και, το κυριότερο, \emph{αραιός}, καθώς κάθε μέτρηση εμπλέκει ελάχιστους
μόνο ζυγούς. Η αραιότητα δεν συνιστά λεπτομέρεια υλοποίησης: είναι εκείνη που καθιστά το
πρόβλημα επιλύσιμο σε πραγματικό χρόνο για δίκτυα χιλιάδων ζυγών. Μια πυκνή παραγοντοποίηση
θα απαιτούσε χρόνο ανάλογο του $n^3$, ενώ η αραιή κλιμακώνεται σχεδόν γραμμικά.

\subsubsection*{Παρατηρησιμότητα}

Η επίλυση της \eqref{eq:gain} προϋποθέτει αντιστρεψιμότητα του $\mathbf{G}$, η οποία δεν
είναι δεδομένη: μεγάλο πλήθος μετρήσεων δεν συνεπάγεται κατ' ανάγκη \emph{κατάλληλες}
μετρήσεις. Εφόσον όλες συγκεντρώνονται σε ένα τμήμα του δικτύου, το υπόλοιπο παραμένει
άγνωστο ανεξαρτήτως πλήθους.

Το σύστημα χαρακτηρίζεται \textbf{παρατηρήσιμο} όταν οι διαθέσιμες μετρήσεις επαρκούν για
τον μονοσήμαντο προσδιορισμό ολόκληρης της κατάστασης. Ο έλεγχος πραγματοποιείται με δύο
ισοδύναμες προσεγγίσεις:

\begin{itemize}
\item \textbf{Τοπολογική}: αναζητείται συνδετικό δένδρο (spanning tree) του δικτύου πλήρως
      καλυμμένο από μετρήσεις. Είναι ταχεία και προσφέρει άμεση φυσική εποπτεία.
\item \textbf{Αριθμητική}: ελέγχεται ο βαθμός του $\mathbf{H}$ ή η παραγοντοποίηση του
      $\mathbf{G}$. Είναι γενικότερη και εντοπίζει και τις οριακές περιπτώσεις.
\end{itemize}

Όταν το σύστημα δεν είναι πλήρως παρατηρήσιμο, εντοπίζονται οι \emph{παρατηρήσιμες νησίδες}
και η κατάσταση εκτιμάται μόνο εντός αυτών. Τα κενά γεφυρώνονται με \emph{ψευδομετρήσεις}
(pseudo-measurements), δηλαδή εκτιμήσεις φορτίου από ιστορικά δεδομένα ή προβλέψεις, στις
οποίες αποδίδεται μεγάλη τυπική απόκλιση ώστε να βαρύνουν ελάχιστα.

\subsubsection*{Ανίχνευση και αναγνώριση εσφαλμένων μετρήσεων}

Στο σημείο αυτό αξιοποιείται ο πλεονασμός. Εσφαλμένη μέτρηση είναι ασύμβατη με τις
υπόλοιπες, και η ασυμβατότητα αυτή αποτυπώνεται στα \emph{υπόλοιπα}
$\mathbf{r} = \mathbf{z}-\mathbf{h}(\hat{\mathbf{x}})$. Η κλασική διαδικασία περιλαμβάνει
δύο στάδια.

\textbf{Στάδιο 1 --- Ανίχνευση.} Υπό την υπόθεση ότι όλες οι μετρήσεις είναι ορθές, η
αντικειμενική συνάρτηση $J(\hat{\mathbf{x}})$ στη λύση ακολουθεί κατανομή $\chi^2$ με
$m-n$ βαθμούς ελευθερίας. Στο προηγούμενο αριθμητικό παράδειγμα, με $m = 400$ και
$n = 199$, οι βαθμοί ελευθερίας είναι 201 και η αναμενόμενη τιμή του $J$ επίσης περίπου
201. Υπέρβαση του κατωφλίου $\chi^2_{m-n,\,\alpha}$ για δεδομένο επίπεδο εμπιστοσύνης
$\alpha$ συνιστά στατιστική ένδειξη ύπαρξης εσφαλμένων δεδομένων. Το τεστ ωστόσο
αποφαίνεται μόνο για την \emph{ύπαρξη} του προβλήματος, όχι για τον εντοπισμό του.

\textbf{Στάδιο 2 --- Αναγνώριση.} Τα ακατέργαστα υπόλοιπα δεν είναι συγκρίσιμα μεταξύ τους,
καθώς κάθε μέτρηση χαρακτηρίζεται από διαφορετική ακρίβεια και διαφορετικό βαθμό στήριξης
από τις γειτονικές. Κανονικοποιούνται επομένως ως

\begin{equation}
r_i^{N} = \frac{|r_i|}{\sqrt{\Omega_{ii}}}, \qquad
\boldsymbol{\Omega} = \mathbf{R} - \mathbf{H}\mathbf{G}^{-1}\mathbf{H}^{T}
\label{eq:normres}
\end{equation}

όπου $\boldsymbol{\Omega}$ ο πίνακας συνδιακύμανσης των υπολοίπων. Η μέτρηση με το μέγιστο
$r_i^{N}$, σύμφωνα με το κριτήριο του μεγίστου κανονικοποιημένου υπολοίπου (Largest
Normalized Residual), θεωρείται ύποπτη, απορρίπτεται, και η εκτίμηση επαναλαμβάνεται χωρίς
αυτήν. Η διαδικασία επαναλαμβάνεται έως ότου το τεστ $\chi^2$ ικανοποιηθεί. Στην πράξη
χρησιμοποιείται συχνά κατώφλι $r_i^{N} > 3$, το οποίο αντιστοιχεί σε γεγονός τριών τυπικών
αποκλίσεων.

\subsubsection*{Δύο θεμελιώδεις περιορισμοί}

Η μέθοδος είναι ισχυρή, υπόκειται όμως σε δύο όρια που πρέπει να είναι σαφή.

\textbf{Πρώτον, οι κρίσιμες μετρήσεις.} Μια μέτρηση χαρακτηρίζεται \emph{κρίσιμη} όταν η
αφαίρεσή της καθιστά το σύστημα μη παρατηρήσιμο. Τέτοιες μετρήσεις εμφανίζουν εξ ορισμού
μηδενικό υπόλοιπο, καθώς δεν υφίσταται άλλη πληροφορία ικανή να τις διαψεύσει και η
εκτίμηση αναγκάζεται να τις αναπαραγάγει επακριβώς. Συνεπώς, σφάλμα σε κρίσιμη μέτρηση
είναι \emph{μη ανιχνεύσιμο}, ανεξαρτήτως μεγέθους. Ο εντοπισμός των κρίσιμων μετρήσεων ενός
δικτύου αποτελεί επομένως μέρος του σχεδιασμού του συστήματος μέτρησης και όχι ακαδημαϊκή
άσκηση.

\textbf{Δεύτερον, οι επιθέσεις έγχυσης ψευδών δεδομένων.} Οι Liu, Ning και Reiter ανέδειξαν
το 2009 ένα ουσιωδώς σοβαρότερο ζήτημα. Έστω επιτιθέμενος με γνώση της τοπολογίας και των
παραμέτρων του δικτύου, δηλαδή του πίνακα $\mathbf{H}$. Εφόσον αλλοιώσει τις μετρήσεις κατά
$\mathbf{a} = \mathbf{H}\mathbf{c}$ για οποιοδήποτε διάνυσμα $\mathbf{c}$, η αλλοιωμένη
μέτρηση $\mathbf{z}_{\text{bad}} = \mathbf{z} + \mathbf{a}$ αντιστοιχεί επακριβώς στην
κατάσταση $\hat{\mathbf{x}} + \mathbf{c}$, ενώ το υπόλοιπο παραμένει \emph{αμετάβλητο}:

\begin{equation}
\mathbf{r}_{\text{bad}} = \mathbf{z} + \mathbf{H}\mathbf{c}
        - \mathbf{H}(\hat{\mathbf{x}} + \mathbf{c}) = \mathbf{r}
\label{eq:fdia}
\end{equation}

Η επίθεση διαφεύγει πλήρως των ανωτέρω στατιστικών ελέγχων, ακριβώς επειδή είναι
\emph{συνεπής} προς το μοντέλο. Δεν πρόκειται για ελάττωμα υλοποίησης, αλλά για δομικό όριο
κάθε ελέγχου που βασίζεται στα υπόλοιπα. Οι επιθέσεις αυτές, γνωστές ως επιθέσεις έγχυσης
ψευδών δεδομένων (False Data Injection Attacks, FDIA), καθώς και οι στρατηγικές
αντιμετώπισής τους, εξετάζονται αναλυτικά στο Κεφάλαιο 8.

\subsubsection*{Ενσωμάτωση μετρήσεων PMU}

Η προσθήκη μετρήσεων PMU μεταβάλλει ουσιωδώς το πρόβλημα. Καθώς η PMU μετρά απευθείας τη
γωνία τάσης, η σχέση μεταξύ μετρήσεων και κατάστασης καθίσταται \emph{γραμμική} όταν
εκφραστεί σε ορθογώνιες συντεταγμένες. Εφόσον το σύστημα είναι πλήρως παρατηρήσιμο μόνο
από PMU, η \eqref{eq:measmodel} εκφυλίζεται σε
$\mathbf{z} = \mathbf{H}\mathbf{x} + \mathbf{e}$ και η λύση δίνεται σε ένα βήμα:

\begin{equation}
\hat{\mathbf{x}} = \left(\mathbf{H}^{T}\mathbf{R}^{-1}\mathbf{H}\right)^{-1}
                    \mathbf{H}^{T}\mathbf{R}^{-1}\mathbf{z}
\label{eq:linse}
\end{equation}

Τα πλεονεκτήματα είναι σημαντικά: δεν απαιτούνται επαναλήψεις, οπότε δεν ανακύπτουν
προβλήματα σύγκλισης ούτε εξάρτηση από την αρχική εκτίμηση· ο $\mathbf{H}$ παραμένει
σταθερός για δεδομένη τοπολογία και συνεπώς προϋπολογίζεται· τέλος, η εκτίμηση μπορεί να
εκτελείται με τον ρυθμό των PMU, δηλαδή δεκάδες φορές ανά δευτερόλεπτο αντί για μία φορά
ανά λεπτό.

Στην πράξη, η πλήρης κάλυψη με PMU σπανίζει για λόγους κόστους, οπότε χρησιμοποιούνται
\textbf{υβριδικοί εκτιμητές} που συνδυάζουν συμβατικές μετρήσεις SCADA με
συγχρονοφασιθέτες. Η βασική δυσκολία έγκειται στην ετερογένεια των ρυθμών, δηλαδή στον
συνδυασμό μετρήσεων που ανανεώνονται ανά 4 s με άλλες που ανανεώνονται ανά 20 ms.

\subsection{Στάδια λειτουργίας: από τον προγραμματισμό στον πραγματικό χρόνο}

Η λειτουργία του συστήματος δεν αποφασίζεται σε μία χρονική στιγμή, αλλά μέσα από
διαδοχικές χρονικές κλίμακες. Η λογική είναι απλή και αξίζει να διατυπωθεί ρητά: \emph{όσο
πλησιάζει η ώρα παράδοσης, τόσο ακριβέστερη καθίσταται η πρόβλεψη και τόσο δαπανηρότερη η
διόρθωση}. Κάθε στάδιο διορθώνει τα σφάλματα του προηγούμενου, με λιγότερο διαθέσιμο χρόνο
και περιορισμένες επιλογές. Στο ευρωπαϊκό μοντέλο αγοράς (Target Model) οι κλίμακες αυτές
αντιστοιχούν σε διακριτές αγορές.

\begin{itemize}
\item \textbf{Μακροπρόθεσμος και προθεσμιακός ορίζοντας} (έτη έως εβδομάδες). Μελέτες
      επάρκειας ισχύος, συντονισμός προγραμμάτων συντήρησης γραμμών και μονάδων, εκτίμηση
      των αναγκών σε εφεδρείες, προθεσμιακά προϊόντα αντιστάθμισης κινδύνου.
\item \textbf{Αγορά Επόμενης Ημέρας} (day-ahead). Αποτελεί το κύριο στάδιο
      προγραμματισμού. Οι συμμετέχοντες υποβάλλουν προσφορές για κάθε αγοραία περίοδο της
      επόμενης ημέρας και η αγορά εκκαθαρίζεται με ενιαία σύζευξη σε ευρωπαϊκό επίπεδο
      (Single Day-Ahead Coupling), ώστε η ενέργεια να ρέει από τις φθηνότερες προς τις
      ακριβότερες περιοχές έως τον κορεσμό της διαθέσιμης διασυνοριακής ικανότητας. Ο
      Διαχειριστής ελέγχει στη συνέχεια την τεχνική εφικτότητα του προκύπτοντος
      προγράμματος.
\item \textbf{Ενδοημερήσια Αγορά} (intra-day). Επιτρέπει τη διόρθωση των θέσεων καθώς
      πλησιάζει η ώρα παράδοσης και βελτιώνονται οι προβλέψεις. Έχει ιδιαίτερη σημασία σε
      συστήματα με υψηλή διείσδυση αιολικής και φωτοβολταϊκής παραγωγής: το σφάλμα
      πρόβλεψης αιολικής παραγωγής μειώνεται δραστικά όσο συντομεύει ο ορίζοντας, οπότε η
      δυνατότητα αναπροσαρμογής λίγες ώρες πριν από την παράδοση αποκτά σημαντική
      οικονομική αξία.
\item \textbf{Πραγματικός χρόνος και εξισορρόπηση} (balancing). Ο Διαχειριστής ενεργοποιεί
      προσφορές ενέργειας εξισορρόπησης, ώστε να καλύψει τις εναπομένουσες αποκλίσεις, να
      αντιμετωπίσει διαταραχές και να διαχειριστεί συμφορήσεις.
\end{itemize}

\subsubsection*{Η αγοραία περίοδος}

Η \emph{αγοραία περίοδος} (Market Time Unit, MTU) είναι η χρονική ανάλυση εκκαθάρισης της
αγοράς. Επί σειρά ετών ήταν η ώρα. Από τις 30 Σεπτεμβρίου 2025, με ημέρα παράδοσης την 1η
Οκτωβρίου 2025, η ενιαία σύζευξη επόμενης ημέρας λειτουργεί με αγοραία περίοδο
\textbf{15 λεπτών} σε όλες τις ευρωπαϊκές ζώνες προσφοράς, εναρμονισμένη με την περίοδο
εκκαθάρισης αποκλίσεων. Η μεταβολή δεν συνιστά τεχνική λεπτομέρεια: μικρότερη αγοραία
περίοδος συνεπάγεται ότι οι απότομες μεταβολές παραγωγής και φορτίου αποτυπώνονται στο ίδιο
το πρόγραμμα, αντί να μετακυλίονται στην εξισορρόπηση, με αποτέλεσμα τη μείωση των αναγκών
σε εφεδρείες.

\subsubsection*{Το ελληνικό πλαίσιο}

Στην Ελλάδα το μοντέλο τέθηκε σε ισχύ την 1η Νοεμβρίου 2020. Η Αγορά Επόμενης Ημέρας και η
Ενδοημερήσια Αγορά λειτουργούν από το Ελληνικό Χρηματιστήριο Ενέργειας (ΕΧΕ), ενώ την
\textbf{Αγορά Εξισορρόπησης} λειτουργεί ο ΑΔΜΗΕ. Ο τελευταίος εκτελεί τη Διαδικασία
Ολοκληρωμένου Προγραμματισμού και ενεργοποιεί σε πραγματικό χρόνο τόσο ενέργεια
εξισορρόπησης όσο και τις εφεδρείες που περιγράφονται στη συνέχεια.

\subsection{Απόκριση συχνότητας και εφεδρείες}

\subsubsection*{Αδράνεια και αρχικός ρυθμός μεταβολής}

Η συχνότητα αποτελεί τον πλέον άμεσο δείκτη του ισοζυγίου ισχύος: παραμένει σταθερή όσο
παραγωγή και ζήτηση εξισορροπούνται και αποκλίνει μόλις εμφανιστεί ανισοζύγιο.

Αμέσως μετά από διαταραχή ισχύος $\Delta P$, τυπικά την απώλεια μεγάλης μονάδας, κανένας
ρυθμιστής δεν έχει προλάβει να αντιδράσει. Το έλλειμμα καλύπτεται από την κινητική ενέργεια
των στρεφόμενων μαζών, οι οποίες επιβραδύνονται. Ο αρχικός ρυθμός μεταβολής της συχνότητας
(Rate of Change of Frequency, ROCOF) προκύπτει από την εξίσωση ταλάντωσης και δίνεται από
τη σχέση

\begin{equation}
\frac{\mathrm{d}f}{\mathrm{d}t} = \frac{f_0 \, \Delta P}{2 H S_n}
\label{eq:rocof}
\end{equation}

όπου $f_0$ η ονομαστική συχνότητα, $H$ η σταθερά αδράνειας του συστήματος σε δευτερόλεπτα
και $S_n$ η συνολική ονομαστική φαινόμενη ισχύς των \emph{συγχρονισμένων} μονάδων.

\emph{Παράδειγμα.} Για σύγχρονη περιοχή με $S_n = 300$ GVA, $H = 5$ s και απώλεια
$\Delta P = 3\,000$ MW, ο αρχικός ρυθμός ανέρχεται σε
$50 \times 3000 / (2 \times 5 \times 300\,000) = 0{,}05$ Hz/s. Εφόσον η ίδια απώλεια
εκδηλωθεί σε ώρα χαμηλού φορτίου, με τη μισή συγχρονισμένη ισχύ, ο ρυθμός διπλασιάζεται.

Η σχέση \eqref{eq:rocof} εξηγεί τον λόγο για τον οποίο η αντικατάσταση σύγχρονων
γεννητριών από μονάδες συνδεδεμένες μέσω μετατροπέων — οι οποίες δεν προσφέρουν εγγενή
αδράνεια — οδηγεί σε ταχύτερες αποκλίσεις συχνότητας. Ο διαθέσιμος χρόνος έως ότου η
συχνότητα φτάσει σε επικίνδυνα επίπεδα συρρικνώνεται και η απόκριση οφείλει να καταστεί
αντιστοίχως ταχύτερη. Από εδώ πηγάζει το ενδιαφέρον για συνθετική αδράνεια και ταχεία
απόκριση συχνότητας.

Το Σχήμα~\ref{fig:freq} δείχνει τη συνολική απόκριση για το συμβάν αναφοράς και
διακρίνει τις τρεις διαδοχικές φάσεις που αναλύονται στη συνέχεια.

\begin{figure}[!htbp]
\centering
\includegraphics[width=\linewidth]{figures/frequency-response.svg}
\caption{Απόκριση συχνότητας σε απώλεια 3\,000 MW. Η κλίση στην αρχή καθορίζεται
αποκλειστικά από την αδράνεια· η εφεδρεία διατήρησης ανακόπτει την πτώση και σταθεροποιεί
τη συχνότητα σε νέα μόνιμη τιμή· η εφεδρεία αποκατάστασης εξαλείφει την εναπομένουσα
απόκλιση και επαναφέρει τα 50 Hz.}
\label{fig:freq}
\end{figure}

\subsubsection*{Εφεδρεία διατήρησης συχνότητας και στατισμός}

Η \textbf{εφεδρεία διατήρησης συχνότητας} (Frequency Containment Reserve, FCR) υλοποιείται
από τους ρυθμιστές στροφών των μονάδων, οι οποίοι μεταβάλλουν την παραγωγή αναλογικά προς
την απόκλιση συχνότητας. Ο συντελεστής αναλογίας εκφράζεται μέσω του \emph{στατισμού} $R$:

\begin{equation}
\frac{\Delta P}{P_n} = -\frac{1}{R}\,\frac{\Delta f}{f_0}
\label{eq:droop}
\end{equation}

Τυπικές τιμές είναι $R = 4$ έως 5\%. \emph{Παράδειγμα:} μονάδα 500 MW με $R = 5\%$, υπό
απόκλιση $\Delta f = -200$ mHz, δηλαδή $-0{,}4\%$, αυξάνει την παραγωγή της κατά
$(0{,}004/0{,}05) \times 500 = 40$ MW.

Το αρνητικό πρόσημο εξασφαλίζει αρνητική ανάδραση, καθώς πτώση της συχνότητας προκαλεί
αύξηση της παραγωγής. Ουσιώδες είναι ωστόσο ότι η ρύθμιση είναι \emph{αναλογική} και όχι
ολοκληρωτική: για να διατηρηθεί η επιπλέον παραγωγή, οφείλει να διατηρηθεί και η απόκλιση
συχνότητας. Η εφεδρεία διατήρησης επομένως \textbf{σταθεροποιεί} τη συχνότητα σε νέα μόνιμη
τιμή, χωρίς να την \textbf{επαναφέρει} στα 50 Hz. Ακριβώς αυτό το εναπομένον σφάλμα
καλείται να εξαλείψει το επόμενο επίπεδο.

\subsubsection*{Εφεδρεία αποκατάστασης συχνότητας σε απομονωμένο σύστημα}

Η \textbf{εφεδρεία αποκατάστασης συχνότητας} (Frequency Restoration Reserve, FRR) εξετάζεται
πρώτα στην απλούστερη περίπτωση: ένα \emph{απομονωμένο} σύστημα, χωρίς διασυνδέσεις με
γειτονικά συστήματα.

Εκεί υπάρχει ένα και μόνο κριτήριο, η συχνότητα. Ο κεντρικός ελεγκτής εφαρμόζει
\emph{ολοκληρωτική} δράση επί της απόκλισης:

\begin{equation}
\Delta P_{\mathrm{FRR}}(t) = -K_I \int_0^{t} \Delta f(\tau)\,\mathrm{d}\tau
\label{eq:frr-isolated}
\end{equation}

Ο ολοκληρωτικός όρος είναι εκείνος που κάνει τη διαφορά. Όσο υπάρχει απόκλιση συχνότητας,
το ολοκλήρωμα συνεχίζει να αυξάνεται και η εντολή παραγωγής να μεγαλώνει· η διαδικασία
σταματά μόνο όταν $\Delta f = 0$. Με άλλα λόγια, η ολοκληρωτική δράση εξαλείφει εξ ορισμού
το μόνιμο σφάλμα που αφήνει η αναλογική δράση της εφεδρείας διατήρησης. Καθώς η συχνότητα
επανέρχεται στα 50 Hz, η συνεισφορά της FCR μηδενίζεται σύμφωνα με την \eqref{eq:droop}:
η εφεδρεία αποκατάστασης \emph{απελευθερώνει} την εφεδρεία διατήρησης, ώστε αυτή να είναι
και πάλι διαθέσιμη για την επόμενη διαταραχή.

Η δράση αυτή είναι σκοπίμως βραδύτερη της FCR. Δύο ελεγκτές που δρουν στο ίδιο μέγεθος με
συγκρίσιμες ταχύτητες αλληλεπιδρούν και μπορούν να οδηγήσουν σε ταλαντώσεις· ο διαχωρισμός
των χρονικών κλιμάκων — δευτερόλεπτα για την FCR, λεπτά για την FRR — είναι εκείνος που
εξασφαλίζει την ευστάθεια της συνολικής ιεραρχίας.

\subsubsection*{Εφεδρεία αποκατάστασης σε σύστημα πολλών περιοχών}

Σε διασυνδεδεμένο σύστημα το πρόβλημα αλλάζει ποιοτικά. Η συχνότητα είναι \emph{κοινή} σε
ολόκληρη τη σύγχρονη περιοχή: μια απώλεια μονάδας στη Γερμανία μειώνει τη συχνότητα και
στην Ελλάδα. Αν κάθε Διαχειριστής ρύθμιζε μόνο με βάση τη συχνότητα, όλοι θα αντιδρούσαν
ταυτόχρονα στην ίδια διαταραχή, ανεξαρτήτως του πού συνέβη. Το αποτέλεσμα θα ήταν
υπερδιόρθωση, ανεξέλεγκτες αποκλίσεις των διασυνοριακών ροών από τα προγράμματα και εν τέλει
οικονομική μεταφορά κόστους από τον υπαίτιο στους υπόλοιπους.

Απαιτείται επομένως δεύτερο κριτήριο, το οποίο να διακρίνει \emph{πού} εκδηλώθηκε το
ανισοζύγιο. Ο ελεγκτής κάθε περιοχής επιδιώκει τον μηδενισμό του \emph{σφάλματος ελέγχου
περιοχής} (Area Control Error):

\begin{equation}
\mathrm{ACE} = \Delta P_{\mathrm{tie}} + K \, \Delta f
\label{eq:ace}
\end{equation}

όπου $\Delta P_{\mathrm{tie}}$ η απόκλιση των διασυνοριακών ανταλλαγών από το πρόγραμμα και
$K$ ο συντελεστής μεροληψίας συχνότητας της περιοχής. Στο απομονωμένο σύστημα ο πρώτος όρος
απλώς δεν υπάρχει και η \eqref{eq:ace} εκφυλίζεται στο κριτήριο της προηγούμενης
υποενότητας.

Ο συνδυασμός των δύο όρων παράγει την επιθυμητή διάκριση, όπως δείχνει το
Σχήμα~\ref{fig:lfc}:

\begin{itemize}
\item \textbf{Διαταραχή εντός της περιοχής.} Η συχνότητα πέφτει ($\Delta f < 0$) και,
      ταυτόχρονα, οι γείτονες στέλνουν βοηθητική ισχύ, οπότε οι εισαγωγές υπερβαίνουν το
      πρόγραμμα ($\Delta P_{\mathrm{tie}} < 0$). Οι δύο όροι έχουν το ίδιο πρόσημο, το ACE
      λαμβάνει μεγάλη απόλυτη τιμή και η περιοχή διορθώνει.
\item \textbf{Διαταραχή εκτός της περιοχής.} Η συχνότητα πέφτει επίσης, αλλά η περιοχή
      \emph{εξάγει} επιπλέον ισχύ προς τον γείτονα ($\Delta P_{\mathrm{tie}} > 0$). Με
      κατάλληλη επιλογή του $K$ οι δύο όροι αλληλοαναιρούνται και το ACE παραμένει σχεδόν
      μηδενικό: η περιοχή έχει ήδη προσφέρει τη συνεισφορά της μέσω της FCR και δεν
      απαιτείται περαιτέρω ενέργεια.
\end{itemize}

Η ιδιότητα αυτή ονομάζεται \emph{μη αλληλεπιδραστικός έλεγχος} (non-interactive control):
κάθε Διαχειριστής αποκαθιστά τα δικά του ανισοζύγια, ενώ η αλληλεγγύη κατά την πρώτη στιγμή
της διαταραχής παραμένει καθολική. Ο συντελεστής $K$ επιλέγεται κατά προσέγγιση ίσος με τη
συνολική ενεργό ισχύ απόκρισης της περιοχής — άθροισμα της συνεισφοράς των ρυθμιστών και
της απόσβεσης του φορτίου — ώστε η αλληλοαναίρεση να είναι όσο το δυνατόν πληρέστερη.

\begin{figure}[!htbp]
\centering
\includegraphics[width=\linewidth]{figures/lfc-areas.svg}
\caption{Εφεδρεία αποκατάστασης σε απομονωμένο και σε διασυνδεδεμένο σύστημα. Στην πρώτη
περίπτωση αρκεί η απόκλιση συχνότητας· στη δεύτερη απαιτείται και ο όρος των διασυνοριακών
ανταλλαγών, ώστε να διορθώνει μόνο η περιοχή στην οποία εκδηλώθηκε η διαταραχή.}
\label{fig:lfc}
\end{figure}

\subsubsection*{Τα τέσσερα επίπεδα εφεδρείας}

\begin{itemize}
\item \textbf{FCR} --- εφεδρεία διατήρησης συχνότητας. Αυτόματη και τοπική, ενεργοποιείται
      από όλες τις μονάδες της σύγχρονης περιοχής από κοινού. Σταθεροποιεί τη συχνότητα.
\item \textbf{aFRR} --- αυτόματη εφεδρεία αποκατάστασης συχνότητας. Αυτόματη και κεντρική,
      οδηγείται από το ACE. Επαναφέρει τη συχνότητα στα 50 Hz και απελευθερώνει την FCR.
\item \textbf{mFRR} --- χειροκίνητη εφεδρεία αποκατάστασης συχνότητας. Ενεργοποιείται με
      εντολή του Διαχειριστή και αντικαθιστά την aFRR σε παρατεταμένα ανισοζύγια.
\item \textbf{RR} --- εφεδρεία αντικατάστασης. Αποκαθιστά τα επίπεδα των προηγούμενων
      εφεδρειών, ώστε το σύστημα να είναι έτοιμο για επόμενη διαταραχή.
\end{itemize}

Στην παλαιότερη βιβλιογραφία και σε μέρος της πρακτικής τα τρία πρώτα επίπεδα απαντούν ως
πρωτεύουσα, δευτερεύουσα και τριτεύουσα ρύθμιση αντιστοίχως. Η ορολογία που χρησιμοποιείται
εδώ είναι εκείνη του SO GL και προτιμάται, καθώς ονομάζει τα επίπεδα κατά τη \emph{λειτουργία}
τους — διατήρηση, αποκατάσταση, αντικατάσταση — και όχι κατά τη σειρά ενεργοποίησης.

Διακρίνεται ένα σαφές μοτίβο: κάθε επίπεδο είναι βραδύτερο, οικονομικότερο και πιο
στοχευμένο από το προηγούμενο, ενώ ο ρόλος του συνίσταται στην \emph{απελευθέρωση} του
ταχύτερου, ώστε αυτό να καταστεί εκ νέου διαθέσιμο.

\begin{table}[!htbp]
\centering
\begin{tabular}{lllc}
\hline
Εφεδρεία & Λειτουργία & Ενεργοποίηση & Πλήρης ενεργοποίηση \\
\hline
FCR  & διατήρηση      & αυτόματη, τοπική   & 30 s \\
aFRR & αποκατάσταση   & αυτόματη, κεντρική & 5 min \\
mFRR & αποκατάσταση   & χειροκίνητη        & 12,5 min \\
RR   & αντικατάσταση  & χειροκίνητη        & $>$ 15 min \\
\hline
\end{tabular}
\caption{Επίπεδα εφεδρείας συχνότητας κατά τον SO GL. Οι χρόνοι πλήρους ενεργοποίησης
αφορούν τη σύγχρονη περιοχή της Ηπειρωτικής Ευρώπης.}
\label{tab:reserves}
\end{table}

\subsubsection*{Ποσοτικές απαιτήσεις για την Ηπειρωτική Ευρώπη}

\begin{itemize}
\item Το \textbf{συμβάν αναφοράς} για τη διαστασιολόγηση της FCR είναι απώλεια 3\,000 MW,
      και αντιστοίχως 3\,000 MW στην αντίθετη κατεύθυνση. Η ελάχιστη διαθέσιμη FCR στη
      σύγχρονη περιοχή ανέρχεται ανά πάσα στιγμή σε 3\,000 MW.
\item Η μέγιστη απόκλιση συχνότητας μόνιμης κατάστασης για το συμβάν αναφοράς είναι
      \textbf{200 mHz}. Από τα δύο αυτά μεγέθη προκύπτει ο συντελεστής αναλογίας μεταξύ FCR
      και απόκλισης συχνότητας: $3\,000\ \text{MW} / 0{,}2\ \text{Hz} = 15\,000$ MW/Hz.
\item Το \textbf{τυπικό εύρος συχνότητας} είναι $\pm 50$ mHz και δεν επιτρέπεται να
      παραβιάζεται για περισσότερα από 15\,000 λεπτά ανά έτος. Η μέγιστη στιγμιαία απόκλιση
      ορίζεται σε 800 mHz.
\item Απόκλιση συχνότητας μεγαλύτερη των 200 mHz οδηγεί σε κατάσταση \emph{έκτακτης
      ανάγκης}.
\end{itemize}

\subsection{Σύνοψη}

Η λειτουργία του ΣΜΗΕ στηρίζεται στον κλειστό βρόχο με τον οποίο ξεκίνησε το κεφάλαιο:

\begin{enumerate}
\item Η \textbf{υποδομή εποπτείας} (SCADA, PMU) παράγει μετρήσεις, οι οποίες είναι
      πλεονάζουσες, θορυβώδεις και εν μέρει ασύγχρονες.
\item Η \textbf{εκτίμηση κατάστασης} τις μετατρέπει σε συνεπή και στατιστικά βέλτιστη
      εικόνα του δικτύου, απορρίπτοντας ταυτόχρονα τα προφανώς εσφαλμένα δεδομένα.
\item Τα \textbf{εργαλεία ανάλυσης ασφάλειας} αξιολογούν την εικόνα αυτή έναντι του
      κριτηρίου N-1 και κατατάσσουν το σύστημα σε μία από τις πέντε καταστάσεις λειτουργίας.
\item Οι \textbf{αποφάσεις} υλοποιούνται μέσω των αγορών και των διατάξεων ελέγχου, σε
      κλίμακες που εκτείνονται από τα δευτερόλεπτα της εφεδρείας διατήρησης έως τους μήνες
      του προγραμματισμού.
\end{enumerate}

Ο βρόχος είναι τόσο ισχυρός όσο ο ασθενέστερος κρίκος του, και ο κρίκος αυτός είναι συχνά
η αξιοπιστία της ίδιας της εικόνας: κρίσιμες μετρήσεις χωρίς πλεονασμό, επιθέσεις συνεπείς
προς το μοντέλο, τμήματα του δικτύου που δεν είναι παρατηρήσιμα.

Οι υπηρεσίες που στηρίζουν τον βρόχο αυτόν ονομάζονται \emph{επικουρικές υπηρεσίες} και
διακρίνονται σε σχετικές και μη σχετικές με τη συχνότητα· τις εξετάζουν αντιστοίχως τα
Κεφάλαια 3 και 2. Το αντίστοιχο πρόβλημα στο επίπεδο της διανομής, όπου η παρατηρησιμότητα
είναι πολύ χαμηλότερη και κυριαρχούν οι ψευδομετρήσεις, εξετάζεται στο Κεφάλαιο 4. Η
ασφάλεια της ίδιας της υποδομής μέτρησης και ελέγχου αποτελεί το αντικείμενο του
Κεφαλαίου 8.

\section{Παραδείγματα}
\label{kef-01-paradeigmata}
% TODO: παραδείγματα και λυμένες ασκήσεις

\section{Βιβλιογραφία}
\label{kef-01-vivliografia}

\subsection*{Βιβλία αναφοράς}

\begin{itemize}
\item A. Abur and A. Gómez Expósito, \textit{Power System State Estimation: Theory and
      Implementation}. New York: Marcel Dekker, 2004.
\item N. G. Hingorani and L. Gyugyi, \textit{Understanding FACTS: Concepts and Technology of
      Flexible AC Transmission Systems}. New York: IEEE Press, 2000.
\item P. Kundur, \textit{Power System Stability and Control}. New York: McGraw-Hill, 1994.
\item A. Monticelli, \textit{State Estimation in Electric Power Systems: A Generalized
      Approach}. Boston, MA: Kluwer Academic Publishers, 1999.
\item A. G. Phadke and J. S. Thorp, \textit{Synchronized Phasor Measurements and Their
      Applications}, 2nd ed. Cham: Springer, 2017.
\item A. J. Wood, B. F. Wollenberg and G. B. Sheblé, \textit{Power Generation, Operation, and
      Control}, 3rd ed. Hoboken, NJ: Wiley, 2014.
\end{itemize}

\subsection*{Θεμελιώδεις εργασίες}

\begin{itemize}
\item T. E. Dy Liacco, «The Adaptive Reliability Control System», \textit{IEEE Transactions
      on Power Apparatus and Systems}, vol. PAS-86, no. 5, pp. 517--531, May 1967.
      doi: \texttt{10.1109/TPAS.1967.291728}
\item Y. Liu, P. Ning and M. K. Reiter, «False Data Injection Attacks against State
      Estimation in Electric Power Grids», \textit{ACM Transactions on Information and System
      Security}, vol. 14, no. 1, art. 13, May 2011.
      doi: \texttt{10.1145/1952982.1952995}
\item A. Monticelli, «Electric Power System State Estimation», \textit{Proceedings of the
      IEEE}, vol. 88, no. 2, pp. 262--282, Feb. 2000.
      doi: \texttt{10.1109/5.824004}
\item F. C. Schweppe and J. Wildes, «Power System Static-State Estimation, Part I: Exact
      Model», \textit{IEEE Transactions on Power Apparatus and Systems}, vol. PAS-89, no. 1,
      pp. 120--125, Jan. 1970. doi: \texttt{10.1109/TPAS.1970.292678}
\item F. C. Schweppe and D. B. Rom, «Power System Static-State Estimation, Part II:
      Approximate Model», \textit{IEEE Transactions on Power Apparatus and Systems},
      vol. PAS-89, no. 1, pp. 125--130, Jan. 1970.
\item F. C. Schweppe, «Power System Static-State Estimation, Part III: Implementation»,
      \textit{IEEE Transactions on Power Apparatus and Systems}, vol. PAS-89, no. 1,
      pp. 130--135, Jan. 1970.
\end{itemize}

\subsection*{Πρότυπα}

\begin{itemize}
\item IEC/IEEE 60255-118-1:2018, \textit{Measuring Relays and Protection Equipment ---
      Part 118-1: Synchrophasor for Power Systems --- Measurements}. IEC/IEEE, 2018.
\item IEC 61850, \textit{Communication Networks and Systems for Power Utility Automation}. IEC.
\item IEC 60870-5-104, \textit{Telecontrol Equipment and Systems --- Part 5-104: Network
      Access for IEC 60870-5-101 Using Standard Transport Profiles}. IEC.
\end{itemize}

\subsection*{Κανονιστικό πλαίσιο και Διαχειριστές}

\begin{itemize}
\item Commission Regulation (EU) 2017/1485 of 2 August 2017 establishing a guideline on
      electricity transmission system operation (System Operation Guideline, SO GL).
      \textit{Official Journal of the European Union}, L 220, 25.8.2017.
      \href{https://eur-lex.europa.eu/eli/reg/2017/1485/oj}{eur-lex.europa.eu/eli/reg/2017/1485/oj}
\item Commission Regulation (EU) 2017/2195 of 23 November 2017 establishing a guideline on
      electricity balancing (Electricity Balancing Guideline, EB GL). \textit{Official Journal
      of the European Union}, L 312, 28.11.2017.
      \href{https://eur-lex.europa.eu/eli/reg/2017/2195/oj}{eur-lex.europa.eu/eli/reg/2017/2195/oj}
\item ENTSO-E, \textit{Synchronous Area Framework Agreement for Regional Group Continental
      Europe --- Policy 1: Load-Frequency Control and Reserves}, 2019.
      \href{https://www.entsoe.eu/}{entsoe.eu}
\item ΑΔΜΗΕ Α.Ε., \textit{Κώδικας Διαχείρισης του Ελληνικού Συστήματος Μεταφοράς Ηλεκτρικής
      Ενέργειας}. \href{https://www.admie.gr/}{admie.gr}
\item ΑΔΜΗΕ Α.Ε., \textit{Κανονισμός Αγοράς Εξισορρόπησης}. \href{https://www.admie.gr/}{admie.gr}
\item Ελληνικό Χρηματιστήριο Ενέργειας (ΕΧΕ), \textit{Κανονισμός Αγοράς Επόμενης Ημέρας και
      Ενδοημερήσιας Αγοράς}. \href{https://www.enexgroup.gr/}{enexgroup.gr}
\end{itemize}
